\section{Introduction}

Social media platforms have reshaped global communication, serving as primary infrastructure for news dissemination, political discourse, and social interaction. This connectivity has also been exploited by \textit{social bots}---automated or semi-automated accounts controlled by algorithmic scripts. The scale of the problem is substantial: a 2025 study analysing approximately 200 million users across seven global events estimates that social media activity comprises roughly 20\% bots, which differ systematically from humans in linguistic patterns, interaction structure, and content dissemination~\cite{ng2025socialbots}. The downstream costs are significant as well; the World Economic Forum ranks disinformation among the top global risks for 2025, and the U.S. Federal Trade Commission reported that scams originating on social media caused \$2.1 billion in losses in 2025, with investment fraud accounting for over \$1.1 billion~\cite{ftc2026scams}. Operating as coordinated networks, these accounts are deployed to manipulate public opinion, amplify disinformation, and distort fake-follower economies. As bots grow more sophisticated, detecting them within complex human-machine interaction graphs has become an important priority across the Artificial Intelligence and Cybersecurity research agendas.

Early detection systems relied primarily on node-level feature engineering, analyzing textual patterns or account metadata. However, these methods quickly became obsolete as bots evolved to hijack legitimate profile semantics. Consequently, recent state-of-the-art architectures have pivoted to Graph Neural Networks (GNNs)~\cite{quang2026integrating}. By framing bot detection as a node classification problem over a heterogeneous interaction graph, modern frameworks successfully integrate local semantic features with multi-relational topologies (e.g., ``follows'', ``retweets'', ``mentions'').

Despite these advances, deployed GNN detection systems face two structural limitations. The first is \textbf{limited robustness to heterophily}. Standard message-passing GNNs largely build on the assumption of \textit{homophily}---the tendency of similar nodes to connect to one another. While early-generation bots formed dense, homophilic ``bot-nets'', many modern bots engage in \textit{relation camouflage}: by deliberately following and interacting with large numbers of genuine human users, they create a highly heterophilic local neighbourhood. During standard GNN aggregation, this camouflage tends to dilute the bot's anomalous signature with the benign features of its human neighbours, a failure mode often described as over-smoothing by camouflage.

The second limitation is an \textbf{explainability gap}. Many advanced GNNs provide limited rationale beyond a binary classification. In real-world platform moderation, administrators are often reluctant to suspend large numbers of accounts based solely on an opaque probability score, given the risk of public backlash and regulatory scrutiny. Moderators typically require concrete, interpretable evidence---such as specific suspicious interaction subgraphs or coordinated community motifs---indicating \textit{why} an account was flagged.

To address the dual challenges of adversarial camouflage and limited transparency, we propose \textbf{XHBot} (eXplainable Heterophily-aware Bot detector). XHBot reconsiders the aggregation process by down-weighting adversarial edges before message passing and maintaining channel separation between homophilic communities and heterophilic camouflage. It further couples detection accuracy with a hierarchical framework for transparent evidence generation.

In summary, the primary contributions of this paper are fourfold:
\begin{itemize}
    \item \textbf{Heterophily-Aware Architecture:} We introduce a Spectral-Guided Topology Refinement (SGTR) module and a Tri-Channel Heterophily-Aware Aggregation (THCA) mechanism that together reduce relation camouflage and mitigate the dilution of bot signals during aggregation.
    \item \textbf{Latent Disentanglement:} We design a Contrastive Prototype Disentanglement (CPD) objective that encourages the model to decouple an account's intrinsic behavioural signature from the noise of its social positioning.
    \item \textbf{Hierarchical Forensic Explainability:} We propose an integrated Hierarchical Forensic Explanation (HFE) module that extracts interpretable evidence at the instance (subgraph) and macro (community-motif) levels to support transparent platform moderation.
    \item \textbf{Strong Empirical Performance:} Empirical evaluations on camouflaged interaction graphs show that XHBot achieves an F1 score of 0.9474 on TwiBot-20, improving over recent GNN baselines by 9.64\% under a unified evaluation protocol, while also providing interpretable evidence for its decisions.
\end{itemize}